\title{\vspace{-1.4cm}\bfseries VedGraph: Agentic Construction of an\\
Evidence-Typed Knowledge Graph for the Vedic Sa\d{m}hit\=as}

\author{%
  Himanshu Mohanty\\[2pt]
  \small Independent researcher\\
  \small \url{https://github.com/HimanshuMohanty-Git24/VedAnvaya}
}
\date{\small Preprint. Measurements frozen at commit \code{6bf220c},
20 September 2026, from a working tree carrying this paper's own untracked
directory and nothing else (\cref{sec:appendix-repro}).}

\maketitle
\thispagestyle{fancy}

\begin{abstract}
\noindent
Large language models can read a historical corpus far faster than a philologist can,
and that speed is the problem: a model proposes fluent claims that are indistinguishable,
once written into a graph, from claims a source actually makes. We report the
construction of \textbf{VedGraph}, a knowledge graph over the four Vedic Sa\d{m}hit\=as
in a single named recension each --- \d{R}gveda \'S\=akala, S\=amaveda Kauthuma
\=Arcika, V\=ajasaneyi M\=adhyandina, and Atharvaveda \'Saunaka, 20{,}210 canonical
mantras in total --- built over seventeen days by a human owner directing bounded
\mbox{LLM} agents alongside deterministic pipelines. The graph holds 165{,}737 nodes and
510{,}905 relationships.

The method rests on two commitments. First, identity is source-independent: a passage's
identifier is a version-5 \mbox{UUID} over a canonical \mbox{URN} built from recension
and structural coordinates alone, and every normalisation --- accent folding, script
transliteration, sandhi-insensitive comparison --- is a comparison instrument that may
generate candidates but may never establish that two passages are the same. Second,
every relationship carries the epistemic layer that produced it: source-explicit,
deterministically derived, model-extracted, or interpretive. Because the layer is a
stored, indexed property rather than documentation, its distribution is measurable.
It is the paper's central empirical result: although agents were used throughout
construction, \textbf{0.17\,\% of relationships are model-extracted} (890 of 510{,}905),
against 21.39\,\% source-explicit and 77.22\,\% deterministically derived. The ratio
depends on the population and we give four; over the semantic assertion layer alone,
where a model is the proposing mechanism, the share is 7.0\,\% --- and none of it was
promoted. The agentic
contribution was concentrated in candidate generation, measurement, cross-file
reconciliation and adversarial refutation --- not in assertion. A policy expressed as
data rather than as convention enforces this: a model may write exactly one lifecycle
status, \code{CANDIDATE}, and the promotion gate is a frozen set of trust classes that
does not contain the model's.

We describe the validation regime that made the ratio hold, in which each layer is
attacked by an independently implemented gate, and we report its failures as evidence
about method. The project deleted 41{,}796 of its own relationships in one revision and
retired 30{,}539 in the next; withdrew 96\,\% of its semantic role fillers rather than
ship a measured 24.8\,\% error; refused an entire corpus on an inter-reader agreement of
0.505; and refused a semantic-resemblance layer after a re-implemented lexical control
outperformed it. In the case we treat at greatest length, a gate built to attack the
project's own \emph{withholdings} --- not only its assertions --- found 39 withheld
S\=amavedic verses whose stated reason the repository's own released data disproved;
the reason was corrected and no verse moved from withheld to released. We state
plainly what no amount of this machinery establishes: of 35{,}131 semantic assertions,
zero have been reviewed by a human, and the project's own certification says so.
We argue that in scholarly knowledge-graph construction the durable contribution of
agents is not the claims they generate but the refutations they can be made to
perform, and that this is only safe when the graph can say, per edge, how each claim
came to be.
\end{abstract}

\vspace{2pt}
\noindent\textbf{Keywords:} knowledge graph construction; \mbox{LLM} agents; provenance;
evidence typing; adversarial validation; digital humanities; computational philology;
Sanskrit; Vedic corpus; text reuse.

\vspace{4pt}
\hrule
\vspace{6pt}
